% ===================================================================
%  तालिका सं. 7 — जाड़े में गाँव। — वसंत में खेत।
%  Traduction 2026 d'apres texto/io, controlee sur texto/fr et
%  texto/en. Consigne : voir texto/hi/10-talika-01.tex.
% ===================================================================

%%K t07-noto-1 noto
\VUnotes{143.2mm}{%
(*) भोजन के ब्योरे के लिए तालिकाएँ 6 और 13 देखिए।%
}

%%K t07-apar-1 apar
\VUsaut{4.0mm}
\VUcentre{13.2pt}{90}{दूसरी शृंखला}
\VUsaut{3.0mm}
\VUfilet{16mm}
\VUsaut{5.6mm}
\VUtitre{15.0pt}{60}{तालिका सं. 7}
\VUsaut{3.0mm}

%%K t07-tit-1 sub
\VUcentre{12.0pt}{0}{\textit{पहला दृश्य।}}
\VUsaut{3.0mm}

%%K t07-tit-2 sub
\VUpk{10.2pt}{40}{जाड़े में गाँव।}
\VUsaut{4.0mm}

%%K t07-c1-01-1 p
1. --- नए वर्ष की छुट्टियों में मैं अपने पिता के साथ नयागाँव गया, एक छोटा गाँव
जहाँ उन्हें कुछ काम निपटाना था।

%%K t07-c1-02-1 p
2. --- हमने वह \VUgras{डाक-गाड़ी} \textsuperscript{(11)} ली जो शहर और उस गाँव के
बीच चलती है; पर ठंड बहुत थी, और इतनी असुविधाजनक गाड़ी में यात्रा तनिक भी सुखद न
रही।

%%K t07-c1-03-1 p
3. --- इसलिए हमें सचमुच ख़ुशी हुई जब हमने \VUgras{गाड़ीवान} को चौक में
\textit{श्वेत भालू} की \VUgras{सराय} \textsuperscript{(2)} के सामने गाड़ी रोकते
देखा। \VUgras{सरायवाला} \textsuperscript{(4)} अपनी देहरी पर हमारी बाट जोह रहा था,
चुपचाप अपना \VUgras{चिलम} पीता हुआ।

%%K t07-c1-03-2 p
उसने हमें भीतर बुलाया, और चूल्हे में चटकती बेल की आग के सामने बैठने के लिए मुझे
दुबारा कहलाने की ज़रूरत न पड़ी। फिर मुझे उस कँपाती \VUgras{उत्तरी हवा} की कोई
चिंता न रही जो पुराने \VUgras{बोर्ड} \textsuperscript{(3)} को चरमराती थी और चिमनी
से निकलते \VUgras{धुएँ} \textsuperscript{(1)} को काले बगूलों में उड़ा ले जाती थी।

%%K t07-c1-04-1 p
4. --- दोपहर के भोजन का समय था। और देर किए बिना हमने उस सादे पर स्वादिष्ट भोजन
के साथ पूरा न्याय किया जो हमें परोसा गया (*)।

%%K t07-tit-3 sub
\VUpk{10.2pt}{40}{कुछ भेंटें।}
\VUsaut{3.0mm}

%%K t07-c1-05-1 p
5. --- भोजन के बाद मैं अपने पिता के साथ गया, जो बीमे के एजेंट हैं, उनके
ग्राहकों से मिलने। पहले हम \VUgras{लोहार} \textsuperscript{(5)} के पास गए, जो
सराय के बगल में रहता है। वह एक घोड़े की नाल ठोंकने में लगा था। मैंने उसके
सहायक के बल की सराहना की, जो घोड़े की \VUgras{टाप} को अपने चमड़े के एप्रन पर
कसकर थामे था, जब तक लोहार भारी हथौड़े की चोटों से \VUgras{नाल}
\textsuperscript{(6)} जड़ता रहा।

%%K t07-c1-06-1 p
6. --- फिर हम गाँव के \VUgras{मोची} \textsuperscript{(7)} बूढ़े जैकब के पास गए।
यद्यपि उसके बोर्ड पर एक शानदार \VUgras{बूट} बना था, वह पुराने घिसे जूतों में नया
तला लगाकर काम चला रहा था।

%%K t07-c1-06-2 p
बूढ़े ने हँसते हुए कहा : « शहर में मैं “जूता बनाने वाला” था और मेरे ग्राहक ही न
थे। यहाँ मैं “मोची” हूँ और काम की कोई कमी नहीं। »

%%K t07-c1-07-1 p
7. --- उसने हमें अपने पड़ोसी \VUgras{हजामतवाले} \textsuperscript{(8)} की बात
बताई, जिसके ग्राहक सब असंतुष्ट रहते हैं। उसके \VUgras{उस्तरे} सदा खंडित रहते हैं
और उसकी \VUgras{कैंची} कभी ठीक नहीं काटती। पर चूँकि वही गाँव का अकेला
हजामतवाला है, लोगों के पास उसी से हजामत बनवाने और बाल कटवाने के सिवा चारा नहीं।
वह ऊपर से कंजूस और रूखा भी है।

%%K t07-c1-08-1 p
8. --- सौभाग्य से दूसरे \VUgras{पड़ोसी} उस जैसे नहीं। \VUgras{गाड़ी बनाने वाला}
\textsuperscript{(9)}, उदाहरण के लिए, भला आदमी है, मेहनती और उपकारी। उसे गाड़ी
मरम्मत को देना सुख की बात है। उसका काम सदा साफ़-सुथरा होता है।

%%K t07-c1-09-1 p
9. --- बूढ़े जैकब को \VUgras{जीनसाज़} \textsuperscript{(10)} से भी कोई शिकायत न
थी, जिसकी वह काम के दबाव में कभी-कभी \VUgras{साज} सीने में मदद करता है।

%%K t07-c1-09-2 p
मेरे पिता ने उससे उसके परिवार के बारे में पूछा। « सब भले-चंगे हैं। मेरी पोती
\VUgras{पाठशाला} \textsuperscript{(12)} में है। उसकी \VUgras{अध्यापिका}
\textsuperscript{(13)} उससे बहुत प्रसन्न हैं; वह अच्छी \VUgras{विद्यार्थिनी}
\textsuperscript{(14)} है। वह मेरे उस शरारती बेटे जैसी नहीं; उसे खेल के सिवा कुछ
नहीं सूझता। \VUgras{प्रधान अध्यापक} \textsuperscript{(15)} उससे कुछ भी नहीं
सिखा पाते। »

%%K t07-tit-4 sub
\VUsaut{3.0mm}
\VUpk{10.2pt}{40}{गाँव की बातें।}
\VUsaut{3.0mm}

%%K t07-c1-10-1 p
10. --- फिर बूढ़े ने हमें गाँव की ख़बरें सुनाईं। एक नई पाठशाला बननी थी, क्योंकि
\VUgras{ग्राम-भवन} वाली बहुत छोटी पड़ गई थी। यह सहज होता, क्योंकि
\VUgras{चक्कीवाले} के बग़ीचे का एक हिस्सा ख़रीद लेना ही बहुत था। बेचारे के लिए
यह बड़ा अच्छा रहता। जब से ठंड पड़ी, \VUgras{चक्की} \textsuperscript{(16)} के
\VUgras{पाट} गेहूँ नहीं पीस सके, क्योंकि \VUgras{पहिया} \textsuperscript{(17)}
\VUgras{बर्फ़} \textsuperscript{(25)} में फँस गया था, जो \VUgras{नाले}
\textsuperscript{(23)} की थी।

%%K t07-c1-11-1 p
11. --- « निर्माण की बात चली तो — हमारा नया \VUgras{गिरजा} \textsuperscript{(18)}
देखा आपने? बर्फ़ की चादर के नीचे वह बड़ा सुरम्य लगता है। \VUgras{कौवे}
\textsuperscript{(19)}, जो अपना पुराना घंटाघर अब पहचानते ही नहीं, उदास होकर
\VUgras{कब्रिस्तान} \textsuperscript{(20)} के बड़े पेड़ पर बैठे रहते हैं। »

%%K t07-c1-12-1 p
12. --- कब्रिस्तान का नाम आते ही मोची को वे लोग याद आ गए जिन्हें मेरे पिता जानते
थे और जो पिछले साल चल बसे। « कितनी \VUgras{कब्रें} \textsuperscript{(21)}, मेरे
भले महाशय, \VUgras{सरो} \textsuperscript{(22)} के नीचे और जुड़ गईं! मौत हमारे
बूढ़े मुंशी को ले गई। सारा गाँव उनकी \VUgras{अंत्येष्टि} में था। »

%%K t07-c1-13-1 p
13. --- « वह रहे उनके उत्तराधिकारी, नाले पर बर्फ़ फिसल रहे हैं। थोड़ी देर पहले
ही आए थे कि मैं उनकी \VUgras{बर्फ़-जूती} \textsuperscript{(26)} का तसमा जोड़ दूँ,
जो टूट गया था। »

%%K t07-c1-14-1 p
14. --- « और वह रहा, नाले के किनारे, नया \VUgras{गाँव का चौकीदार}
\textsuperscript{(27)}। वह पुराना सिपाही है और गाँव के शरारती लड़कों का काल। उसके
\VUgras{पट्टे} \textsuperscript{(29)} और उसकी \VUgras{झालरदार टोपी}
\textsuperscript{(28)} दिखते ही वे भाग खड़े होते हैं। »

%%K t07-c1-15-1 p
15. --- « अरे! वह रहा बूढ़ा जोसेफ़, नानबाई के लिए \VUgras{लकड़ी के गट्ठरों की
गाड़ी} \textsuperscript{(30)} लाता हुआ। --- ठीक कहा, मेरे पिता बोले, मैं उसके
\VUgras{ख़च्चरों} \textsuperscript{(31)} की जोड़ी पहचानता हूँ। मुझे उससे एक बात
करनी है, यही मौका है। नमस्ते, बूढ़े जैकब। »

%%K t07-c1-16-1 p
16. --- हम बाहर निकल ही रहे थे कि गाँव के बच्चे पाठशाला से छूटकर
\VUgras{बर्फ़ की फिसलपट्टी} \textsuperscript{(33)} पर टूट पड़े, इस ख़तरे के साथ
कि \VUgras{गिरने} \textsuperscript{(34)} पर हाथ-पाँव तोड़ बैठें। उनमें से एक ने
गाड़ीवाले के यहाँ से अपनी \VUgras{स्लेज} \textsuperscript{(32)} लाई, उस पर अपने
एक साथी को बिठाया, और दौड़ पड़ा।

%%K t07-c1-17-1 p
17. --- दूसरे एक \VUgras{बर्फ़ के ढेर} \textsuperscript{(37)} की ओर लपके, जो
\VUgras{पुल} \textsuperscript{(40)} के पास था, और उस \VUgras{बर्फ़ के आदमी} \textsuperscript{(39)}
पर मारने लगे जो उन्होंने कल बनाया था, जिसे उन्होंने एक टोपी, एक चिलम और कंधे पर
लटका एक बस्ता दे रखा था।

%%K t07-c1-18-1 p
18. --- उन्हें बर्फ़ीली हवा और ठंड की परवाह न थी। उनमें से अधिकतर के पास
\VUgras{गुलूबंद} \textsuperscript{(35)} तक न था, और \VUgras{बर्फ़ के गोलों}
\textsuperscript{(38)} की लड़ाई के बाद गरम होने के लिए मुट्ठी भर तपते
\VUgras{चेस्टनट} \textsuperscript{(42)} ही उन्हें बहुत थे, जो वे बूढ़ी जैकलीन
\VUgras{बेचनेवाली} से ख़रीदते थे, जो अपनी बहुत पतली \VUgras{चादर}
\textsuperscript{(43)} में ठंड से काँपती रहती है।

%%K t07-tit-5 sub
\VUsaut{3.0mm}
\VUcentre{12.0pt}{0}{\textit{दूसरा दृश्य।}}
\VUsaut{3.0mm}

%%K t07-tit-6 sub
\VUpk{10.2pt}{40}{वसंत में खेत का घर।}
\VUsaut{4.0mm}

%%K t07-c2-01-1 p
1. --- जाड़े के सब सुखों के बावजूद \VUgras{वसंत} के लौटने का हम गहरे आनंद से
स्वागत करते हैं।

%%K t07-c2-01-2 p
मई के महीने में शाम को खेत का आँगन कैसा उत्साह भरा चित्र बनाता है!

%%K t07-c2-02-1 p
2. --- क्षितिज पर \VUgras{सूरज} \textsuperscript{(1)} धीरे-धीरे ढल रहा है, और
\VUgras{देहात} \textsuperscript{(3)} को अपनी लाल \VUgras{किरणों} \textsuperscript{(2)}
से भर रहा है। मेहनती \VUgras{हलवाहा} \textsuperscript{(6)} अपना \VUgras{खेत}
\textsuperscript{(4)} जोतने की जल्दी में है।

%%K t07-c2-02-2 p
अपने \VUgras{हल} \textsuperscript{(5)} से, जिसे एक मज़बूत \VUgras{घोड़ा}
\textsuperscript{(7)} खींचता है, वह \VUgras{कूँड़} \textsuperscript{(8)} खींचता
है, जिसमें
\VUgras{बोने वाला} \textsuperscript{(9)} मुट्ठी-मुट्ठी वह \VUgras{बीज} छिड़केगा
जिससे सुनहरी बालियाँ निकलेंगी।

%%K t07-c2-03-1 p
3. --- अपनी हँसियों से \VUgras{घास काटने वालों} \textsuperscript{(11)} ने मैदान की
घास काट ली है, सुखाने वालों ने \VUgras{सूखी घास} \textsuperscript{(10)} को
\VUgras{पंजों} \textsuperscript{(12)} और खुरपों से उलट-पलटकर सुखा लिया है, और अब
वे लौट रहे हैं।

%%K t07-c2-03-2 p
कुछ बैलों से खींची भारी गाड़ी के आगे-आगे चलते हैं, कंधे पर अपने औज़ार लिए।
दूसरे, अधिक आलसी, सुगंधित सूखी घास पर आराम से पसर गए हैं।

%%K t07-c2-04-1 p
4. --- चलते-चलते वे \VUgras{माली} \textsuperscript{(16)} का मित्रता से अभिवादन
करते हैं, जो \VUgras{बग़ीचे} \textsuperscript{(13)} में लगा है, \VUgras{फलदार
पेड़ों} की छँटाई करता और \VUgras{नाशपाती के पेड़ों} \textsuperscript{(14)} में
कलम लगाता। \VUgras{आलूबुख़ारे के पेड़} \textsuperscript{(15)} फूलों से लदे हैं, और
अच्छी खाद पाए \VUgras{सब्ज़ी के बग़ीचे} \textsuperscript{(17)} में सब्ज़ियाँ,
गोभी और सलाद पहले से अच्छी हो रही हैं।

%%K t07-c2-04-2 p
नीची दीवार पर एक सुंदर \VUgras{मोर} \textsuperscript{(18)} अपने शानदार पंख दिखाने
को पूँछ फैलाता है।

%%K t07-tit-7 sub
\VUpk{10.2pt}{40}{खेत का आँगन।}
\VUsaut{3.0mm}

%%K t07-c2-05-1 p
5. --- \VUgras{अस्तबल} \textsuperscript{(19)} के दरवाज़े खुले हैं, और चरनी तथा
\VUgras{बिछावन की पुआल} \textsuperscript{(23)} दिखाई देती है उस \VUgras{घोड़ी}
\textsuperscript{(21)} की, जिसे \VUgras{साईस} \textsuperscript{(20)} ने अभी
खोला है। \VUgras{बछेड़ा} \textsuperscript{(22)}, अपनी लंबी टाँगों पर इतना
मज़ेदार, कुलाँचें भरता अपनी माँ के पीछे चलता है।

%%K t07-c2-06-1 p
6. --- \VUgras{कुएँ} \textsuperscript{(29)} के पास \VUgras{गायें}
\textsuperscript{(25)} उस लकड़ी की \VUgras{नाँद} \textsuperscript{(26)} से पानी
पीने जाती हैं जो उनका घाट है। \VUgras{बछड़ा} \textsuperscript{(24)} इसी मौके पर
दूध पी लेता है, और \VUgras{बैल} \textsuperscript{(27)}, चारे की अच्छी गंध पाकर,
प्रसन्नता से रँभाता है और अपने बलवान \VUgras{सींगों} \textsuperscript{(28)} वाला
सिर गर्व से उठाता है।

%%K t07-c2-07-1 p
7. --- सब काम में लगे हैं। \VUgras{नौकरानी} \textsuperscript{(32)} कुएँ की
\VUgras{रस्सी} \textsuperscript{(31)} थामे है। वह जानवरों को पिलाने के लिए एक
\VUgras{बालटी} \textsuperscript{(30)} में पानी खींच रही है। \VUgras{खलिहान}
\textsuperscript{(33)} के पास एक आदमी बिखरी पुआल बटोर रहा है, और
\VUgras{खेत के घर} \textsuperscript{(34)} के दरवाज़े के सामने एक बूढ़ी
\VUgras{कातने वाली} \textsuperscript{(35)}, अपने चरखे और अपनी \VUgras{तकली} के
साथ, पुराने दिनों की तरह \VUgras{सन} या \VUgras{पटसन} कात रही है।

%%K t07-tit-8 sub
\VUsaut{3.0mm}
\VUpk{10.2pt}{40}{खेत का मालिक।}
\VUsaut{3.0mm}

%%K t07-c2-08-1 p
8. --- \VUgras{खेत का मालिक} \textsuperscript{(36)} यह सारा काम चलाता है और अपने
आदमियों को हुक्म देता है। वह उनसे कहता है कि \VUgras{हेंगा} \textsuperscript{(40)}
और \VUgras{कुदाल} \textsuperscript{(39)} ढककर रख दें, जो उस \VUgras{घर}
\textsuperscript{(38)} के पास छोड़ दिए गए हैं जो \VUgras{कुत्ते}
\textsuperscript{(37)} का है।

%%K t07-c2-09-1 p
9. --- उसकी हाँक से एक \VUgras{कबूतर} \textsuperscript{(41)} चौंक उठता है, जो
पूरी गति से \VUgras{कबूतरख़ाने} \textsuperscript{(42)} की ओर उड़ जाता है, और
\VUgras{मुर्ग़ियाँ} \textsuperscript{(43)} भी, जो जल्दी-जल्दी
\VUgras{दड़बे} \textsuperscript{(44)} में लौट जाती हैं।

%%K t07-c2-10-1 p
10. --- \VUgras{भेड़शाला} \textsuperscript{(48)} के सामने, यह रहा बूढ़ा
\VUgras{गड़रिया} \textsuperscript{(45)} अपना \VUgras{रेवड़} \textsuperscript{(49)}
लौटाता हुआ। अपनी \VUgras{लाठी} \textsuperscript{(46)} थामे, जो उसका साथ उतना ही
नहीं छोड़ती जितना उसका फीका \VUgras{कुरता} \textsuperscript{(47)}, वह बाड़े की ओर
\VUgras{मेढ़े} \textsuperscript{(50)}, \VUgras{भेड़ें} \textsuperscript{(53)},
\VUgras{मेमने} \textsuperscript{(54)}, \VUgras{बकरियाँ} \textsuperscript{(51)} और
\VUgras{बकरी के बच्चे} \textsuperscript{(52)} हाँकता है। उसका वफ़ादार
\VUgras{कुत्ता} \textsuperscript{(55)} आर्गुस सबसे पीछे आता है और भौंककर पिछड़ों
को आगे करता है।

%%K t07-c2-11-1 p
11. --- यह सारा शोर-शराबा उस भली \VUgras{दुधारू गाय} \textsuperscript{(56)} की
शांति नहीं बिगाड़ता, जो अपनी \VUgras{घंटी} \textsuperscript{(57)} बजाती रहती है, जब उसे
\VUgras{गोशाला} \textsuperscript{(58)} के दरवाज़े के सामने दुहा जाता है।

%%K t07-c2-12-1 p
12. --- वह जैसे समझती है कि उसे अपना दूध देना ही है, ताकि
\VUgras{खेत की मालकिन} \textsuperscript{(60)} \VUgras{मथानी} \textsuperscript{(61)}
में \VUgras{मक्खन} बना सके, या वे बढ़िया \VUgras{पनीर} \textsuperscript{(64)} जो
\VUgras{टब} \textsuperscript{(63)} पर रखे तख़्ते पर निथर रहे हैं।

%%K t07-c2-13-1 p
13. --- सारी \VUgras{दुग्धशाला} \textsuperscript{(59)} \VUgras{खेत के मज़दूर}
\textsuperscript{(65)} की बदौलत बहुत साफ़ रहती है, जिसने अभी उस बड़ी बालटी में
\VUgras{दूध के डिब्बे} \textsuperscript{(62)} धोए हैं जो वह उठाए है।

%%K t07-tit-9 sub
\VUsaut{3.0mm}
\VUpk{10.2pt}{40}{मुर्ग़ीख़ाना।}
\VUsaut{3.0mm}

%%K t07-c2-14-1 p
14. --- अपनी भारी खड़ाऊँ में वह कुछ भद्दे ढंग से चलता है, और इससे
\VUgras{टर्की} \textsuperscript{(68)}, \VUgras{बत्तखें} \textsuperscript{(66)},
\VUgras{बत्तख़ों के नर} \textsuperscript{(67)} और \VUgras{हंस} \textsuperscript{(69)}
उड़ जाते हैं, जो अपनी \VUgras{चोंचें} \textsuperscript{(70)} चौड़ी खोलकर बेचैन
कोलाहल मचाते हैं।

%%K t07-c2-14-2 p
\VUgras{मुर्ग़ा} \textsuperscript{(71)}, जो अधिक लड़ाकू है, अपनी लाल
\VUgras{कलगी} \textsuperscript{(72)} उठाता है, अपनी \VUgras{पूँछ}
\textsuperscript{(73)} के पंख झाड़ता है और एक गूँजती « कुकड़ूँ-कूँ » छोड़ता है।

%%K t07-c2-15-1 p
15. --- एक \VUgras{मुर्ग़ी-माँ} \textsuperscript{(74)}
\VUgras{गोबर के ढेर} \textsuperscript{(81)} पर चुग रही है, अपने \VUgras{चूज़ों}
\textsuperscript{(75)} के साथ। \VUgras{सूअर का बच्चा} \textsuperscript{(78)} अपनी माँ की ओर दौड़ जाता है,
उस विशाल \VUgras{सूअरी} \textsuperscript{(77)} की ओर जो हमें बाड़े के पास दिखाई
देती है।

%%K t07-c2-16-1 p
16. --- अपने दड़बों में \VUgras{ख़रगोश} \textsuperscript{(79)} वह कोमल
\VUgras{घास} \textsuperscript{(80)} चाव से खा रहे हैं जो मालिक की बेटी उन्हें लाती
है। जेनी को अपने ख़रगोश बहुत प्रिय हैं, और हर दिन, दड़बे से ताज़े \VUgras{अंडे}
\textsuperscript{(76)} लाने के बाद, वह अपने नन्हे मित्रों से मिलने आती है।
\VUsaut{6.0mm}
\VUfilet{22mm}
